% =====================================================================
% Rewritten Introduction, Research Questions, and Literature Review
% Scoped to RQ1 only. RQ2 (B-REASO transfer) and RQ3 (confidence-based
% stopping) move to Future Work.
%
% CITATION KEYS: existing keys reused where they already appear in your
% .bib. New keys are marked % NEW -- add these to references.bib.
% =====================================================================

\section{Introduction}

Large language models can now solve competition-level mathematics, but nearly
all of that progress has happened in English. Bengali has more than 230 million
speakers and is still treated as low-resource. It is thinly represented in
pretraining corpora, and reasoning-specific Bengali data barely exists. Whether
a model can actually reason \emph{in} Bengali, rather than translate its way
through a problem, is still an open question.

Two recent efforts approach this gap from opposite directions. TigerLLM
\parencite{raihan2025tigerllm} continually pretrained a 1B-parameter model on
roughly 10M tokens of Grade 6--12 Bengali textbook content, plus instruction
data distilled from frontier models. It reportedly outperformed much larger
generalist systems, which the authors read as evidence that targeted,
high-quality pretraining can substitute for scale. TituLLM
\parencite{nahin2025titullm} took the opposite route, pretraining on roughly 37B
tokens of general-domain Bengali with an extended tokenizer, and found gains
that were capability-specific rather than uniform. Read side by side the two
results look contradictory. They are not really a test of the same question,
though, because in both cases model size and pretraining design move together.
Neither study isolates which factor is doing the work.

That is the first gap this paper addresses. A second follows from it. If
reasoning ability can be elicited rather than taught, how much data does that
take? The less-is-more literature argues the answer is surprisingly little. LIMA
\parencite{zhou2023lima} % NEW
proposed the Superficial Alignment Hypothesis, which holds that a model's
knowledge is acquired almost entirely during pretraining and that alignment
mainly teaches it which format to answer in. It demonstrated competitive
instruction-following from 1,000 curated examples. LIMO \parencite{ye2025limo}
extended the claim to mathematical reasoning, and s1
\parencite{muennighoff2025s1} reproduced it with 1,000 traces plus
inference-time budget forcing. Both LIMO and s1, however, demonstrate their
result at 32B parameters and in English. Whether minimal-data elicitation
survives the move to a genuinely low-resource language, and to the model sizes
most researchers can actually train, has not been tested.

We test both questions with a single controlled experiment. Seven backbones
spanning 0.6B to 4B parameters are fine-tuned with LoRA on an identical set of
curated Bengali mathematical reasoning traces, then evaluated on a held-out
Bengali mathematics benchmark before and after training. The backbones are
chosen so that pretraining depth and parameter scale vary separately.
TigerLLM-1B and TituLLM-1B match in size but differ by roughly $3{,}700\times$
in targeted Bengali pretraining. TituLLM-1B and TituLLM-3B share a pretraining
recipe and differ in size. The Qwen3 family (0.6B, 1.7B, 4B) provides a scale
ladder with essentially no Bengali-specific pretraining at all.

Our contributions are:

\begin{itemize}
  \item \textbf{A controlled separation of scale from targeted pretraining
        depth in Bengali.} Within a single model family holding pretraining
        constant, accuracy rises from 3.6\% at 0.6B to 16.8\% at 1.7B to 44.8\%
        at 4B, and every adjacent step on that ladder is significant (Fisher
        exact, $p < 10^{-5}$). Targeted Bengali pretraining produces no
        comparable gradient: four Bengali-specialised models spanning 1B to
        3.2B parameters all sit between 0.4\% and 4.4\%.
  \item \textbf{No comparison in this study favours the model with more
        Bengali pretraining.} TigerLLM-1B and TituLLM-1B are matched in size
        and differ by roughly $3{,}700\times$ in targeted Bengali pretraining,
        and the model with less of it scores higher, 4.4\% against 2.0\%. The
        sharpest case is TituLLM-3B-Instruct against Qwen3-1.7B: roughly twice
        the parameters plus a Bengali continued-pretraining corpus, against
        neither, and it loses 1.2\% to 16.8\% ($p < 10^{-4}$). This contradicts
        the substitution claim in \parencite{raihan2025tigerllm}, tested on a
        Bengali benchmark using that work's own released model.
  \item \textbf{A scale boundary on minimal-data elicitation.} LIMO-style
        fine-tuning produced no significant accuracy change on any of seven
        backbones. That includes one model whose 44.8\% baseline shows LIMO's
        stated precondition of latent capability is clearly satisfied.
  \item \textbf{Evidence that elicitation changes form without changing
        capability.} No backbone retained even half the items it had answered
        correctly before tuning, and Qwen3-4B abandoned its own answer format
        on 222 of 250 items while its accuracy moved by 13. Mean Bengali
        character adherence fell in six of seven models, by up to 17 points,
        and the per-item decline is significant in each of those six (exact
        sign test, $p < 10^{-7}$).
  \item \textbf{A direct test of the difficulty-calibration explanation.}
        Retraining the same backbone on difficulty-matched easy data performed
        significantly worse than the hard-data adapter, 12.4\% against 18.0\%
        (exact McNemar, $p = 0.039$), which rules out the most obvious
        alternative account of the null result.
  \item \textbf{A measurement of the Bengali tokenizer tax} and what it means
        for which models can be trained on consumer hardware.
\end{itemize}

\subsection{Research Questions}

\textbf{RQ1:} When a model is given a small set of Bengali mathematical examples
intended to elicit reasoning, what primarily determines its performance:
parameter scale, or the amount of Bengali seen during pretraining? We answer
this by comparing models that differ in only one of these factors at a time.

\textbf{RQ1a:} Does minimal-data elicitation, demonstrated at 32B parameters in
English, transfer to the 0.6B to 4B range in Bengali?

\textbf{RQ1b:} If elicitation does not improve accuracy, does it change model
behaviour at all, and in which direction?

TigerLLM-1B and TituLLM-1B match in size but differ in pretraining depth.
TituLLM-1B and TituLLM-3B share a pretraining approach but differ in size. The
Qwen3 ladder isolates scale with Bengali exposure held near zero. This design
lets us attribute any difference to one factor at a time, which no prior work in
this space has done.


% =====================================================================
\section{Literature Review}

\subsection{The Low-Resource Reasoning Baseline}

Reasoning in low-resource languages is a known weak point, and Bengali is no
exception. \textcite{alnazi2025eval} tested six LLMs on four Bengali tasks under
zero-shot, few-shot and chain-of-thought prompting. The best result, GPT-4o with
three-shot CoT, reached 99\% F1 on sentiment, but CoT was not reliably helpful
and in several cases reduced accuracy through hallucination. The failure mode is
therefore not simply insufficient capacity. Reasoning behaviour itself is
unstable in this setting.

BenNumEval \parencite{ahmed2025bennumeval} sharpens the point. Across 3.2k
samples in six numerical reasoning categories,
DeepSeek-R1-Distill-Llama-70B, built from top-tier English reasoning
distillation, collapsed to 4.7\% accuracy in native Bengali. That is far below
Gemini-2.0-flash's 79.6\%, and below several non-distilled baselines. This is
not merely a failure to transfer; it suggests distilling English reasoning
traces can actively damage existing multilingual ability. SOMADHAN
\parencite{paul2025somadhan} adds an important qualifier. Using 8.8k
expert-translated Bengali word problems, Llama-3.3-70B reached 88\%, but smaller
models largely failed to adapt to native-language reasoning at all. The
elicitation problem therefore interacts with model size, which is precisely what
this paper sets out to measure.

\subsection{The Backbone Debate: Scale versus Targeted Depth}

This is the central disagreement in the literature, and it splits three ways
rather than two.

The \textbf{targeted-depth} position is TigerLLM \parencite{raihan2025tigerllm}.
Pretrained on only about 10M tokens of Bengali textbook content plus 100k
instruction pairs distilled from GPT-4o and Claude, its 1B variant reportedly
outperformed 11B generalist models, reaching 68\% on BanglaRQA. The claim is not
that small models can be competitive. It is that targeted depth substitutes for
scale outright.

The \textbf{scale} position pushes back. \textcite{yong2025crosslingual} found
that test-time scaling becomes reliable only at 3B parameters or above, with a
14B model at a sweet spot, and noted that Bengali's tokenizer makes reasoning
roughly $3.5\times$ more expensive than English. MathMist
\parencite{sobhani2026mathmist}, a 30k-pair benchmark across 13 languages,
reports that scaling Qwen from 0.6B to 14B shrank the Bengali--English accuracy
gap from 7\% to 1\%, with smaller models showing what the authors call
linguistic brittleness during CoT.

Classical scaling work supplies the mechanism. \textcite{kaplan2020scaling} % NEW
showed that language modelling loss depends strongly on parameter count and only
weakly on architectural shape, to the point where aspect ratio can vary by a
factor of 40 with little effect. \textcite{wei2022emergent} % NEW
is more directly relevant to our range. They report that three-digit arithmetic
emerges only at roughly 13B parameters, and that chain-of-thought prompting
surpasses standard prompting on GSM8K only at around 100B parameters. If
arithmetic competence and CoT benefit both emerge above 13B, then the 0.6B to 4B
range studied here sits entirely below the threshold at which the relevant
abilities have been shown to appear.

A third position complicates both. TituLLM \parencite{nahin2025titullm}
pretrained on roughly 37B tokens of general-domain Bengali with an extended
tokenizer, and found that this helped commonsense-style tasks but not
broad-knowledge tasks such as MMLU. The authors partly attribute this to
single-epoch training and the absence of instruction tuning. On this account
pretraining depth helps some capabilities and not others.

The problem common to all three is that none of them isolates scale from
pretraining depth. TigerLLM's depth claim exists only at 1B and 9B. TituLLM's
scale claim exists only inside its own general-corpus recipe. The two variables
are always entangled. Closing that gap, by fixing scale at 1B while varying
depth and fixing the recipe while varying scale, is the design of this paper.

\subsection{The Less-Is-More Elicitation Paradigm}

A second strand argues that surfacing reasoning in an already-capable model
takes remarkably little data. This is what makes a small Bengali training set a
reasonable intervention rather than an underpowered shortcut.

LIMA \parencite{zhou2023lima} % NEW
states the underlying claim most directly, as the Superficial Alignment
Hypothesis: a model's knowledge is acquired almost entirely during pretraining,
and alignment principally teaches it which format and style to use when
responding. LIMO \parencite{ye2025limo} carries the argument into mathematics,
reporting that 800 curated reasoning traces elicited more sophisticated
reasoning than training on the 100k-sample NuminaMath dataset, with gains
saturating around 1.6k samples or 32B parameters. s1
\parencite{muennighoff2025s1} arrives somewhere similar by a different route:
1k traces plus budget forcing let s1-32B beat o1-preview on AIME24 by 27\%. LESS
\parencite{xia2024less} offers a mechanism, showing that selecting 5\% of a
270k-point pool by projected low-rank gradient relevance beat training on the
full set. Functional relevance beats raw volume.

Two features of this literature matter here. First, every demonstration is at
32B parameters or above. Second, all of them are in high-resource languages.
LIMO's own stated precondition, that the base model must already hold dense and
relevant knowledge from pretraining, is an empirical claim about the backbone.
Whether it holds at 0.6B to 4B in Bengali is exactly what we test.

\subsection{The Elicitation Tax: Native Reasoning versus English Pivot}

There is a measurable cost to reasoning \emph{in} Bengali specifically.
\textcite{qi2025xreasoning}, using the 370-question, 11-language XReasoning
benchmark, found that forcing native Bengali reasoning instead of pivoting
through English dropped accuracy on hard tasks such as AIME from 26\% to 17\%.
Their proposed fix was tested on at most 817 samples, so it is a partial answer
at best. \textcite{yong2025crosslingual} adds that Bengali's tokenizer makes
reasoning roughly $3.5\times$ more expensive, and that even where in-domain
scaling works, generalisation to out-of-domain cultural knowledge stays weak.
Taken together with the BenNumEval collapse, these results suggest that
interventions meant to improve Bengali reasoning may trade language adherence
against accuracy. We observe that tension directly in our own results.

\subsection{Adaptation Method and Its Limits}

Because this study fine-tunes with LoRA rather than full parameter updates, the
method itself is a variable that has to be accounted for.

LoRA \parencite{hu2021lora} % NEW
freezes pretrained weights and injects trainable low-rank matrices. It was
introduced with results showing parity or better against full fine-tuning on
GLUE-style tasks, for instance 91.3 versus 91.1 average GLUE for DeBERTa-XXL
using 4.7M of 1.5B parameters. Those tasks are largely classification and
generation within a model's existing competence, though.
\textcite{biderman2024lora} % NEW
tests the harder case directly. On mathematics, measured as GSM8K strict match
after continued pretraining, full fine-tuning reached 0.293 while LoRA at rank
256 reached 0.202 and rank 16 only 0.158. On code instruction tuning, full
fine-tuning reached 0.470 against 0.358 for rank 16. Their conclusion is that
``in the standard low-rank settings, LoRA substantially underperforms full
finetuning,'' which they attribute to full fine-tuning finding perturbations of
rank 10 to 100 times greater than typical LoRA configurations. They also report
the converse: LoRA forgets less, keeping the tuned model's behaviour closer to
the base model.

This has two consequences here. It predicts that a rank-16 configuration will
under-acquire new mathematical capability, which bears directly on how our null
result should be read. It also predicts that base behaviour will be
comparatively preserved, which is consistent with what we observe. We note that
Biderman et al. train only at 7B, outside the range studied here.

Domain-adaptive pretraining supplies the other half of the frame. Both TigerLLM
and TituLLM are instances of continued pretraining on a target domain, differing
by more than three orders of magnitude in token budget, and neither reports an
ablation isolating that budget from the base model it starts from.

\subsection{The Cost of Bengali Tokenization}

Tokenizer efficiency is usually treated as an implementation detail. For Bengali
it is a first-order constraint. \textcite{ahia2023alllanguages} % NEW
show that languages with their own scripts require up to $5\times$ more tokens
than English to convey identical content, that API costs rise correspondingly,
and that high fragmentation can stop even a single in-context example from
fitting in a fixed context window. \textcite{petrov2023language} % NEW
quantify this per model, reporting Bengali tokenization premiums that range from
9.65 for GPT-2 and 5.84 for ChatGPT down to 1.01 for the Indic-focused MuRIL.
That is close to a tenfold spread in how much Bengali text fits in the same
budget.

The backbones studied here span that spread, though by accident rather than
design. Llama~3 \parencite{llama3herd} % NEW
adopts a 128k-token vocabulary explicitly to improve non-English compression,
and TituLLM extends it further for Bengali. Gemma~3 \parencite{gemma3} % NEW
ships a 262k-entry SentencePiece vocabulary described as ``more balanced for
non-English languages,'' which TigerLLM inherits without doing any tokenizer
work of its own. Qwen3 \parencite{qwen3} % NEW
uses a 151,669-entry BBPE vocabulary covering 119 languages. These choices were
made for reasons unrelated to Bengali, yet they determine how much Bengali
reasoning fits inside a fixed training budget, and therefore which models can be
trained at all on consumer hardware.

\subsection{Gap Statement}

Three gaps remain open for RQ1.

First, no Bengali-language study separates parameter scale from pretraining
depth as independent factors. The works in Section~2.2 always vary size and
corpus design together, so existing work cannot tell us which one is
responsible.

Second, minimal-data elicitation has been demonstrated only at 32B parameters or
above, and only in high-resource languages, while the abilities it aims to
surface, arithmetic competence and chain-of-thought benefit, have themselves
been reported to emerge only above roughly 13B \parencite{wei2022emergent}. % NEW
Whether the paradigm holds below that threshold, in a language carrying a
$3.5\times$ tokenization penalty, is untested.

Third, the interaction between adaptation method and this question is
unexamined. LoRA is the only practically available option at consumer scale, yet
it is documented to under-acquire new capability relative to full fine-tuning
\parencite{biderman2024lora}, % NEW
and no prior work reports what minimal-data elicitation does under that
constraint in a low-resource language.


% =====================================================================
% NEW BIBTEX ENTRIES -- verify every field against the actual papers
% =====================================================================
%
% @article{kaplan2020scaling,
%   title={Scaling Laws for Neural Language Models},
%   author={Kaplan, Jared and McCandlish, Sam and Henighan, Tom and Brown, Tom B.
%           and Chess, Benjamin and Child, Rewon and Gray, Scott and Radford, Alec
%           and Wu, Jeffrey and Amodei, Dario},
%   journal={arXiv preprint arXiv:2001.08361}, year={2020}}
%
% @article{wei2022emergent,
%   title={Emergent Abilities of Large Language Models},
%   author={Wei, Jason and Tay, Yi and Bommasani, Rishi and Raffel, Colin and
%           Zoph, Barret and Borgeaud, Sebastian and Yogatama, Dani and Bosma, Maarten
%           and Zhou, Denny and Metzler, Donald and Chi, Ed H. and Hashimoto, Tatsunori
%           and Vinyals, Oriol and Liang, Percy and Dean, Jeff and Fedus, William},
%   journal={Transactions on Machine Learning Research}, year={2022}}
%
% @article{hu2021lora,
%   title={LoRA: Low-Rank Adaptation of Large Language Models},
%   author={Hu, Edward J and Shen, Yelong and Wallis, Phillip and Allen-Zhu, Zeyuan
%           and Li, Yuanzhi and Wang, Shean and Wang, Lu and Chen, Weizhu},
%   journal={arXiv preprint arXiv:2106.09685}, year={2021}}
%
% @article{biderman2024lora,
%   title={LoRA Learns Less and Forgets Less},
%   author={Biderman, Dan and Portes, Jacob and Ortiz, Jose Javier Gonzalez and
%           Paul, Mansheej and Greengard, Philip and Jennings, Connor and King, Daniel
%           and Havens, Sam and Chiley, Vitaliy and Frankle, Jonathan and
%           Blakeney, Cody and Cunningham, John P.},
%   journal={Transactions on Machine Learning Research}, year={2024}}
%
% @article{ahia2023alllanguages,
%   title={Do All Languages Cost the Same? Tokenization in the Era of Commercial Language Models},
%   author={Ahia, Orevaoghene and Kumar, Sachin and Gonen, Hila and Kasai, Jungo and
%           Mortensen, David R. and Smith, Noah A. and Tsvetkov, Yulia},
%   journal={arXiv preprint arXiv:2305.13707}, year={2023}}
%
% @article{petrov2023language,
%   title={Language Model Tokenizers Introduce Unfairness Between Languages},
%   author={Petrov, Aleksandar and La Malfa, Emanuele and Torr, Philip H.S. and Bibi, Adel},
%   journal={arXiv preprint arXiv:2305.15425}, year={2023}}   % VERIFY arXiv id
%
% @inproceedings{zhou2023lima,
%   title={LIMA: Less Is More for Alignment},
%   author={Zhou, Chunting and Liu, Pengfei and Xu, Puxin and Iyer, Srini and Sun, Jiao
%           and Mao, Yuning and Ma, Xuezhe and Efrat, Avia and Yu, Ping and Yu, Lili
%           and Zhang, Susan and Ghosh, Gargi and Lewis, Mike and Zettlemoyer, Luke
%           and Levy, Omer},
%   booktitle={Advances in Neural Information Processing Systems}, year={2023}}
%
% @article{llama3herd,
%   title={The Llama 3 Herd of Models},
%   author={{Llama Team, AI@Meta}}, year={2024},
%   journal={arXiv preprint arXiv:2407.21783}}   % VERIFY
%
% @article{gemma3,
%   title={Gemma 3 Technical Report},
%   author={{Gemma Team, Google DeepMind}},
%   journal={arXiv preprint arXiv:2503.19786}, year={2025}}
%
% @article{qwen3,
%   title={Qwen3 Technical Report},
%   author={{Qwen Team}}, year={2025},
%   journal={arXiv preprint arXiv:2505.09388}}   % VERIFY
